% concatenated from SENDOV.tex
%\documentclass[12pt]{article}
\documentclass{amsart}
%	options include 12pt or 11pt or 10pt
%	classes include article, report, book, letter, thesis
\usepackage{setspace}
\usepackage{a4}
\usepackage{amssymb,amsmath,amsthm,latexsym}
\usepackage{amsfonts}
\usepackage{amsfonts}
\usepackage{graphicx}
\usepackage{textcomp}
\usepackage{cite}
\usepackage{enumerate}
\usepackage[mathscr]{euscript}
\usepackage{mathtools}
\newtheorem{theorem}{Theorem}[section]
\newtheorem{acknowledgement}[theorem]{Acknowledgement}
\newtheorem{conjecture}[theorem]{Conjecture}
\newtheorem{corollary}[theorem] {Corollary}
\newtheorem{definition}[theorem]{Definition}
\newtheorem{example}[theorem]{Example}
\newtheorem{lemma} [theorem]{Lemma}
\newtheorem{notation}[theorem]{Notation}
\newtheorem{observation}[theorem]{Observation}
\newtheorem{problem}[theorem]{Problem}
\newtheorem{proposition}[theorem]{Proposition}
\newtheorem{remark}[theorem]{Remark}
\newtheorem{question}[theorem]{Question}
\setlength{\parindent}{0pt} \setlength{\evensidemargin}{0.3cm}
\setlength{\oddsidemargin}{0.3cm} \setlength{\topmargin}{-2cm}
\textwidth 16cm \textheight 23cm
\onehalfspacing
\title{This is the title}
\usepackage{amssymb}
\usepackage{amssymb}
\usepackage{amssymb}
\usepackage{amssymb}
\usepackage{amsmath}
\usepackage{tikz}
\usepackage{hyperref}
\usepackage{enumerate}
\usepackage{mathtools}
\usepackage{amsmath}
\usepackage{tikz}
\usepackage{amssymb}
\usepackage{amsmath}
\usepackage{tikz}
\usepackage{hyperref}
\newcommand{\Cross}{\mathbin{\tikz [x=1.4ex,y=1.4ex,line width=.2ex] \draw (0,0) -- (1,1) (0,1) -- (1,0);}}%
\raggedbottom

\newcommand{\abs}[1]{\lvert#1\rvert}

%    Blank box placeholder for figures (to avoid requiring any
%    particular graphics capabilities for printing this document).
\newcommand{\blankbox}[2]{%
	\parbox{\columnwidth}{\centering
		%    Set fboxsep to 0 so that the actual size of the box will match the
		%    given measurements more closely.
		\setlength{\fboxsep}{0pt}%
		\fbox{\raisebox{0pt}[#2]{\hspace{#1}}}%
	}%
}
\usepackage{fancyhdr}
%\pagestyle{fancy}
%\lhead{C*-ALGEBRAIC SCHUR PRODUCT}
%\rhead{K. MAHESH KRISHNA}
\pagestyle{fancy}
%\thispagestyle{empty}
\fancyhead[LO]{C*-ALGEBRAIC GAUSS-LUCAS THEOREM AND C*-ALGEBRAIC SENDOV'S CONJECTURE}
\fancyhead[RE]{K. MAHESH KRISHNA}



\begin{document}
	\begin{center}
	{\bf{C*-ALGEBRAIC GAUSS-LUCAS THEOREM AND C*-ALGEBRAIC SENDOV'S CONJECTURE}}\\
	{\bf{K. MAHESH KRISHNA}}\\
	Post Doctoral Fellow\\
	Statistics and Mathematics Unit\\
	Indian Statistical Institute, Bangalore Centre\\
	Karnataka 560 059 India\\
	Email: kmaheshak@gmail.com \\
	\today
\end{center}

\hrule
\vspace{0.5cm}


%--------------------------------------
\textbf{Abstract}:  Using a result of Robertson \textit{[Proc. Edinburgh Math. Soc. (2), 1976]}, we introduce a  notion of differentiation of maps on certain classes of unital commutative  C*-algebras. We then derive C*-algebraic Gauss-Lucas theorem and formulate C*-algebraic Sendov's conjecture. We verify C*-algebraic Sendov's conjecture for polynomials of degree 2.  

\textbf{Keywords}: Sendov's conjecture, C*-algebra.

\textbf{Mathematics Subject Classification (2020)}: 30C15, 46L05.
\tableofcontents
%--------------------------------------
\section{Introduction}
Let $a$ be a complex number, $r>0$ be a real number and $\overline{\mathbb{D}(a, r)}$ be the closed disc in $\mathbb{C}$ centered at $a$ of radius $r$. Let $\mathbb{C}[z]$ be the set of all polynomials over $\mathbb{C}$. In 1958, Sendov made the following conjecture which became known as Sendov's conjecture.
 \begin{conjecture}\label{SENDOV} \cite{HAYMAN, MARDEN, RAHMANSCHMEISSERBOOK}	\textbf{(Sendov's conjecture)
 	Let $n \in \mathbb{N}\setminus \{1\}$ and $p(z)=(z-a_1)(z-a_2)\cdots (z-a_n)\in \mathbb{C}[z]$ be such that $a_1, a_2, \dots,  a_n \in \overline{\mathbb{D}(0, 1)}$. Then for each $a_j$, $1\leq j\leq n$, there exists  a zero $b$ of $p'$ such that $b \in  \overline{\mathbb{D}(a_j, 1)}$.} 
 \end{conjecture} 
   A direct calculation says that Conjecture \ref{SENDOV} holds for degree two polynomials. For $n=3$, Conjecture  \ref{SENDOV}  is proved in \cite{BRANNAN}. For $n=3,4$,  Conjecture \ref{SENDOV}  is proved in \cite{RUBINSTEIN}.   For $n=5$,  Conjecture \ref{SENDOV}  is proved in \cite{MEIRSHARMA}. For $n=6$,  Conjecture \ref{SENDOV}  is proved in \cite{KATSOPRINAKIS, BORCEA1996}. For $n=7$,  Conjecture \ref{SENDOV}  is proved in \cite{BROWN1997}. For $n=8$,  Conjecture \ref{SENDOV}  is proved in \cite{BROWNXIANG}. Conjecture \ref{SENDOV}  is proved for all polynomials after certain (unknown) degree in \cite{TAO}. On the other hand, Conjecture \ref{SENDOV}  is proved for various special cases of polynomials and position of roots (inside the closed unit disc) in \cite{BORCEA, BROWN1991, BROWN, BOJANOV, SCHMEISSER, DEGOT, CHALEBGWA, RUBINSTEIN, GOODMANRAHMANRATTI, MILLER, CHIJIWA, KASMALKAR, JOYAL, VAIJAITUZAHARESCU, MILLER2005, MILLER2008, SAFFTWOMEY, MILLER1990, BORCEA1998, RAHMANTARIQ, GOODMAN, PAWLOWSKI, MILLER1992, CHIJIWA2010, SCHMIEDERSZYNAL, DAEPPGORKINVOSS, BOJANOVRAHMANSZYNAL, BYRNE, BYRNE1997}. Non-Archimedean version of Conjecture \ref{SENDOV} has been solved in  \cite{CHOILEE}. In this paper, we derive Gauss-Lucas theorem for polynomials over certain C*-algebras. We then formulate Sendov's conjecture for polynomials over C*-algebras.
 
 
   
   
  
  


\section{C*-algebraic Gauss-Lucas Theorem and C*-algebraic Sendov's Conjecture}
Let $\mathcal{A}$	be a unital commutative C*-algebra and let $G(\mathcal{A})$ be the set of all invertible elements of $\mathcal{A}$. Let $\mathcal{A}[z]$ be the set of all polynomials over $\mathcal{A}$. We assume $G(\mathcal{A})$ is dense in $\mathcal{A}$. Such C*-algebras exist and a characterization of such C*-algebras is given by Robertson  \cite{ROBERTSON}. We now introduce the following definition.
\begin{definition}\label{CDIFFE}
\textbf{(C*-algebraic differentiation) Let $\mathcal{A}$	be a unital commutative C*-algebra and $G(\mathcal{A})$ be  dense in $\mathcal{A}$. Let $f: \mathcal{A}\to \mathcal{A}$ be a function and $\omega\in \mathcal{A}$. We say that	$f$ is C*-algebraic differentiable at $\omega$ if there exists an $L\in \mathcal{A}$ satisfying the following:  for each $\varepsilon>0$, there exists a $\delta>0$ such that if $z\in \mathcal{A}$ satisfies $\|z-\omega\|<\delta$ and $z-\omega\in G(\mathcal{A})$, then 
\begin{align*}
	\|(z-\omega)^{-1}(f(z)-f(\omega))-L\|<\varepsilon.
\end{align*}}
In this case, we write $f'(\omega)=L$.
\end{definition}
We note that, because $G(\mathcal{A})$ is  dense in $\mathcal{A}$, Definition \ref{CDIFFE} is well-defined. Now we derive C*-algebraic Gauss-Lucas theorem. 
\begin{theorem}
\textbf{(C*-algebraic Gauss-Lucas Theorem)} Let $p(z)=(z-a_1)(z-a_2)\cdots (z-a_n)\in\mathcal{A}[z]$. 	Define 
\begin{align*}
	\mathcal{B}\coloneqq \{z \in \mathcal{A}: z-a_j \text{ is  invertible in }  \mathcal{A} \text{ for all } 1\leq j \leq n\}, \quad 	\mathcal{C}\coloneqq \{z \in \mathcal{A}: p'(z)=0\}.	
\end{align*}
 Then for each $z \in \mathcal{B}\cap  \mathcal{C}$, there are positive $\omega_{z_1}, \dots, \omega_{z_n} \in  \mathcal{A}$ such that 
\begin{align}\label{CONVEX}
	z=\sum_{j=1}^{n}\omega_{z_j}a_j, \quad \sum_{j=1}^{n}\omega_{z_j}=1.
\end{align}
\end{theorem}
\begin{proof}
	We have 
	\begin{align*}
		p'(z)=\sum_{j=1}^{n}(z-a_1)\cdots \widehat{(z-a_j)}\cdots (z-a_n), \quad \forall z \in \mathcal{A}.
	\end{align*}
	where the term with cap is missing. Then 
	\begin{align*}
		p(z)^{-1}	p'(z)=\sum_{j=1}^{n}(z-a_j)^{-1}, \quad \forall z \in \mathcal{B}.
	\end{align*}
We therefore have 
\begin{align}\label{FIRST}
	0=\sum_{j=1}^{n}(z-a_j)^{-1}=\sum_{j=1}^{n}((z-a_j)(z-a_j)^*)^{-1}(z^*-a_j^*), \quad \forall z \in \mathcal{B}\cap  \mathcal{C}.
\end{align}
Now  define 
\begin{align*}
	\omega_z \coloneqq \sum_{j=1}^{n}((z-a_j)(z-a_j)^*)^{-1}, \quad \forall z \in \mathcal{B}\cap  \mathcal{C}.
\end{align*}
Note that whenever $ z \in \mathcal{B}$, each $((z-a_j)(z-a_j)^*)^{-1}$ is positive and invertible, $1\leq j \leq n$. Hence $\omega_z$ is positive and invertible. Equation (\ref{FIRST}) then gives 
\begin{align*}
z=\omega_z^{-1}\sum_{j=1}^{n}((z-a_j)(z-a_j)^*)^{-1}a_j, \quad \forall z \in \mathcal{B}\cap  \mathcal{C}.
\end{align*}
\end{proof}
\begin{definition}\label{CDISC}
	\textbf{Given a unital  C*-algebra $\mathcal{A}$ with identity $1$  and an element $a\in \mathcal{A}$, we define the C*-algebraic closed unit disc centered at $a$ and of radius $r>0$, $r\in \mathbb{R}$, denoted as $\overline{\mathbb{D}^*(a, r)}$  by
	\begin{align*}
	 \overline{\mathbb{D}^*(a, r)}\coloneqq 	\{z\in \mathcal{A}: (z-a)(z-a)^*\leq \sqrt{r}\cdot 1\}.
	\end{align*}}
\end{definition}
With Definition \ref{CDISC} we can formulate Conjecture \ref{SENDOV} for C*-algebras.
\begin{conjecture}\label{CSENDOV} 	\textbf{(Commutative C*-algebraic Sendov's conjecture)
Let $\mathcal{A}$	be a unital commutative C*-algebra and $G(\mathcal{A})$ be  dense in $\mathcal{A}$. 	Let $n \in \mathbb{N}\setminus \{1\}$ and $p(z)=(z-a_1)(z-a_2)\cdots (z-a_n)\in\mathcal{A}[z]$ be such that $a_1, a_2, \dots,  a_n \in \overline{\mathbb{D^*}(0, 1)}$. Assume that $p'$ admits roots in $\mathcal{A}$, say $b_1, b_2, \dots,  b_{n-1} \in \overline{\mathbb{D^*}(0, 1)}$ and each $b_k$ can be written in the form of Equation (\ref{CONVEX}). Then for each $a_j$, $1\leq j\leq n$, there exists  a zero $b$ of $p'$ such that $b \in  \overline{\mathbb{D}^*(a_j, 1)}$.} 	
\end{conjecture}
For arbitary C*-algebras, without introducing differentiation we can formulate Conjecture Conjecture \ref{CSENDOV}   as follows. 
\begin{conjecture}\label{CSENDOV2} 
\textbf{(C*-algebraic Sendov's conjecture)
	Let $\mathcal{A}$	be a unital  C*-algebra. 	Let $n \in \mathbb{N}\setminus \{1\}$ and $p(z)=(z-a_1)(z-a_2)\cdots (z-a_n)\in\mathcal{A}[z]$ be such that $a_1, a_2, \dots,  a_n \in \overline{\mathbb{D^*}(0, 1)}$.  Define 
	\begin{align*}
		p'(z)=\sum_{j=1}^{n}(z-a_1)\cdots \widehat{(z-a_j)}\cdots (z-a_n), \quad \forall z \in \mathcal{A}.
	\end{align*}
	where the term with cap is missing. Assume that $p'$ admits roots in $\mathcal{A}$, say $b_1, b_2, \dots,  b_{n-1} \in \overline{\mathbb{D^*}(0, 1)}$ and each $b_k$ can be written in the form of Equation (\ref{CONVEX}). Then for each $a_j$, $1\leq j\leq n$, there exists  a zero $b$ of $p'$ such that $b \in  \overline{\mathbb{D}^*(a_j, 1)}$.} 	
\end{conjecture}
%We can formulate the following variant of Conjecture \ref{CSENDOV2}. 
%\begin{conjecture}
%\textbf{(C*-algebraic Sendov's conjecture - 2)
%	Let $\mathcal{A}$	be a unital  C*-algebra.	Let $n \in \mathbb{N}\setminus \{1\}$ and $p(z)=(z-a_1)(z-a_2)\cdots (z-a_n)\in\mathcal{A}[z]$ be such that $a_1, a_2, \dots,  a_n \in \overline{\mathbb{D^*}(0, 1)}$. Assume that $p'$ completely factors into linear factors. Then for each $a_j$, $1\leq j\leq n$, there exists  a zero $b$ of $p'$ such that $b \in  \overline{\mathbb{D}^*(a_j, 1)}$.} 	
%\end{conjecture}
We end with the following result.
\begin{theorem}
Let $\mathcal{A}$	be a unital  C*-algebra. Conjecture \ref{CSENDOV2}   holds for polynomials  of degree 2 of the form $p(z)=(z-a)(z-b)\in \mathcal{A}[z]$, where $a, b \in \overline{\mathbb{D}^*(0,1)}$. 
\end{theorem}
\begin{proof}
We have 
\begin{align*}
	p'(z)=2z-(a+b), \quad \forall z \in \mathcal{A}.
\end{align*}
Therefore the zero of $p'$ is $\frac{a+b}{2}$ which is in the form of Equation (\ref{CONVEX}). Now we have
\begin{align*}
	2(aa^*+bb^*)\geq (a+b)(a+b)^*, \quad 2(aa^*+bb^*)\geq (a-b)(a-b)^* 
\end{align*}
which gives
\begin{align*}
\left(\frac{a+b}{2}-a\right)\left(\frac{a+b}{2}-a\right)^*\leq 1, \quad \left(\frac{a+b}{2}-b\right)\left(\frac{a+b}{2}-b\right)^*\leq 1.	
\end{align*}
Therefore 
\begin{align*}
\frac{a+b}{2}\in \overline{\mathbb{D}^*\left(a, 1\right)}, \quad \frac{a+b}{2}\in \overline{\mathbb{D}^*(b, 1)}	
\end{align*}
which completes the proof. 
\end{proof}






 
  
  
  
  
  
  
  
 \bibliographystyle{plain}
 \bibliography{reference.bib}
 


\end{document}
